\documentclass{article}

\usepackage{hyperref}

\newcommand{\meetingdate}{20 May 2026}
\usepackage[
    a4paper,
    top=3.5cm,
    bottom=2.3cm,
    left=2.3cm,
    right=2.3cm,
    headheight=42pt
]{geometry}
\usepackage[utf8]{inputenc}
\usepackage[T1]{fontenc}
\usepackage{fancyhdr}
\usepackage{lastpage}
\usepackage{enumitem}
\usepackage{amssymb}
\usepackage{etoolbox}
\usepackage{xcolor}
\usepackage{tikz}
\usepackage{tcolorbox}

\newcommand{\project}{Interactive VR Museum}
\newcommand{\meetingtype}{Weekly Meeting}
\providecommand{\meetingdate}{\today}

\newcommand{\authorname}{Jannick Seper}
\newcommand{\tutors}{Rahel Arnold, Raphael Waltenspül}

\lhead{\project\\\meetingtype\\Date: \meetingdate}
\rhead{\authorname\\\tutors\\Page \thepage\ of \pageref{LastPage}}
\pagestyle{fancy}

\newcommand{\done}{
    \section*{1. What I have done}
}

\newcommand{\support}{
    \section*{2. Where I need support}
}

\newcommand{\plan}{
    \section*{3. What I am planning to do for the next meeting}
}

\newtcolorbox{checkpointbox}{
    colback=gray!5,
    colframe=gray!50,
    boxrule=0.5pt,
    arc=3mm,
    left=3mm,
    right=3mm,
    top=2mm,
    bottom=2mm,
    width=\textwidth
}

\newcommand{\checkpoint}[2][open]{
    \begin{checkpointbox}
        \ifstrequal{#1}{done}
            {$\boxtimes$}
            {$\square$}
        \hspace{0.3cm} #2
    \end{checkpointbox}
}

\begin{document}

\done
\begin{itemize}
    \item Installed Blender, Godot with .NET support, Rider, and Insomnia.
    \item Set up the vitrivr engine first using Docker and later directly on my machine. During this process, I updated several configuration files to match the newer format.
    \item Tested Insomnia with POST requests and successfully received ``retrievables`` from the API.
    \item Started working through the \href{https://docs.godotengine.org/en/stable/getting_started/first_3d_game/index.html}{Official Godot Tutorial}. After completing the chapter ``Jumping and squashing monsters'', the project no longer started correctly. Therefore, I repeated the tutorial using the example project provided on \href{https://github.com/godotengine/godot-demo-projects/tree/master/3d/squash_the_creeps}{GitHub}
    \item Integration and Debugging of XR/OpenXR on Meta Quest 3
    \item Prototype progress:
    \begin{itemize}
        \item Extracted the core ideas of the XRTools virtual keyboard and simplified its implementation to make appearance customization easier.
        \item Created a scene where the user can type a word or text. After pressing Enter, a request is sent to a Flask Python server at: \texttt{http://<serverIP>:5050/search\_one}\\
        The Flask server then forwards the query to the vitrivr engine:\\ \texttt{http://<serverIP>:7070/api/sandbox/query}\\
        The Flask server returns a JSON response containing a URL pointing to the image:\\ \texttt{http://<serverIP>:5050/media/<image>.jpg}\\
        The application then makes a second request to retrieve the image itself. The image is then displayed on a 2D mesh inside a 3D environment.
        \item The frame automatically adjusts its size to the image dimensions. The scaling is currently not correct, so the frame thickness can differ between the horizontal and vertical sides.
        \item All objects (except the ground) can currently be picked up. The keyboard and output image are hidden depending on user input.
        \item Created a second scene containing a simple menu where the Flask server endpoint can be changed(Platform-Switching with double trigger, disabling player movements in menu).
    \end{itemize}
    
    \item Provided my public SSH key and a sketch of the current architecture. The architecture is ``work in progress`` because I have not yet been able to test nginx (see last point in the next section).
\end{itemize}

\support
\begin{itemize}
    \item Access to BlenderKit and Polyhaven (paid products for now?).
    \item Streaming VR content to the Meta Quest from my MacBook is currently not possible because Meta Horizon Link is not available on macOS.
    \item Access to a Vive Focus 3/XR Elite headset
    \item Can I connect a Meta Quest or other VR devices to eduroam?
    \item Regarding vitrivr:\\
    At the moment, the vitrivr engine does not provide the real image path or filename for a retrievable. The VR application therefore only could receive the retrievable ID, but not the actual image (location). Because of this, I have not yet tested nginx because in my understanding nginx can only expose files that already exist in accessible directories via URLs. So I still need the Python/Flask server to map retrievable IDs to the actual image files? Is there a better way to obtain the original image path or filename from a retrievable?
\end{itemize}

\newpage
\plan
\checkpoint[done]{Fix the frame scaling so that all borders have equal thickness}
\checkpoint[done]{Deploy prototype on the Vive devices}
\checkpoint[done]{Try running vitrivr on the server}
\checkpoint{Cleanup Code}
\checkpoint{Create a simple museum environment (room, lighting, etc.)}
\checkpoint{Prepare for the next weekly meeting}
\checkpoint[done]{(Reduce the image retrieval process to a single request)}
\checkpoint[done]{(Test nginx integration)}

\end{document}